\section{Aplikacja mobilna MeshTracker}

\subsection*{Architektura}

Aplikacja wykorzystuje architekturę opartą na wzorcu \textbf{MVVM
(Model-View-ViewModel)}\footnote{\url{https://learn.microsoft.com/pl-pl/dotnet/architecture/maui/mvvm}}
z osobnymi, dodatkowymi warstwami odpowiedzialnymi za komunikację zewnętrzną z aplikacją Meshtastic.
Jest to standardowo wykorzystywany wzorzec, który oddziela logikę działania aplikacji od interfejsu
użytkownika. Warstwa \textit{View} odpowiada za renderowanie interfejsu i reagowanie na zdarzenia
użytkownika, warstwa \textit{ViewModel} przechowuje stan ekranu i przetwarza dane z modelu,
natomiast warstwa \textit{Model} reprezentuje dane domenowe - węzły, strefy geofencingu
i zdarzenia przekroczenia granic. Komunikacja z siecią Meshtastic odbywa się przez dedykowaną
warstwę serwisową, która izoluje resztę aplikacji od szczegółów protokołu AIDL (rysunek~\ref{fig:architektura}).

\begin{figure}[H]
    \label{fig:architektura}
    \centering
    \includegraphics[width=0.8\linewidth]{ArchitectureWhiteBackground.png}
    \caption{Architektura aplikacji MeshTracker}
\end{figure}

\subsection*{Przegląd ekranów aplikacji}

Interfejs użytkownika aplikacji MeshTracker zorganizowany jest wokół trzech głównych ekranów. Wszystkie
dostępne są przez pasek nawigacyjny w dolnej części ekranu: \textbf{mapy}, \textbf{listy węzłów} oraz
\textbf{ustawień}. 

Każdy ekran posiada dedykowany ViewModel odpowiedzialny
za logikę prezentacji, zgodnie z przyjętym wzorcem MVVM. Komponenty współdzielone między
ekranami, takie jak pasek stanu połączenia, wydzielono do osobnego pakietu \texttt{components}.
Wygląd aplikacji definiowany jest przez pakiet \texttt{theme}, zawierający paletę kolorów,
typografię oraz motyw Material Design 3. Strukturę pakietu \texttt{ui} ilustruje rysunek~\ref{fig:ui-structure}.

\begin{figure}[H]
    \centering
    \scriptsize
    \dirtree{%
        .1 ui/.
        .2 MainScreen.kt.
        .2 components/.
        .3 ConnectionStatusBar.kt.
        .2 map/.
        .3 MapScreen.kt.
        .3 MapViewModel.kt.
        .2 nodes/.
        .3 NodeDetailScreen.kt.
        .3 NodeFilterState.kt.
        .3 NodeItem.kt.
        .3 NodeListScreen.kt.
        .2 settings/.
        .3 SettingsScreen.kt.
        .3 SettingsViewModel.kt.
        .2 theme/.
        .3 Color.kt.
        .3 Theme.kt.
        .3 Type.kt.
        .2 zones/.
        .3 ZoneBottomSheet.kt.
        .3 ZoneConfirmDialog.kt.
        .3 ZoneDetailScreen.kt.
        .3 ZoneViewModel.kt.
    }
    \caption{Struktura pakietu \texttt{ui}}
    \label{fig:ui-structure}
\end{figure}

\paragraph{Mapy.}
\texttt{MapScreen} jest głównym ekranem aplikacji. Wyświetla mapę Google Maps,
na której prezentowane są markery węzłów oraz pozycja użytkownika. Każdy marker renderowany
jest jako bitmapa z inicjałami węzła. Kolor markera zależy od roli węzła oraz stanu połączenia.
W prawym górnym rogu markera znajduje się znacznik jakości sygnału. W lewym dolnym rogu ekranu
umieszczono \textit{floating button} otwierający menu geofencingu (rysunek~\ref{fig:map-screen}).

\begin{figure}[!h]
    \centering
    \begin{minipage}{0.48\linewidth}
        \centering
        \includegraphics[width=\linewidth]{mapaZblizenie.jpg}
        \captionof{figure}{Ekran mapy aplikacji MeshTracker}
        \label{fig:map-screen}
    \end{minipage}
    \hfill
    \begin{minipage}{0.48\linewidth}
        \centering
        \includegraphics[width=\linewidth]{wezlyLista.jpg}
        \captionof{figure}{Ekran listy węzłów}
        \label{fig:wezly-lista}
    \end{minipage}
\end{figure}


\paragraph{Markery.}

Wizualizacja markerów została zaprojektowana tak, aby w jak najbardziej przejrzysty sposób
prezentować możliwie wiele danych jednocześnie. Kolor markera odpowiada przypisanej mu roli,
a dodatkowa warstwa kolorów reprezentuje stan połączenia węzła z siecią Meshtastic
(rysunek~\ref{fig:markery-kolory}).

\begin{figure}[!h]
    \centering
    \begin{subfigure}{0.48\linewidth}
        \centering
        \includegraphics[width=\linewidth]{markeryRole.png}
        \caption{Mapowanie roli węzła do koloru}
        \label{fig:markery-role}
    \end{subfigure}
    \hfill
    \begin{subfigure}{0.48\linewidth}
        \centering
        \includegraphics[width=\linewidth]{markeryStanPolaczenia.png}
        \caption{Mapowanie stanu połączenia do koloru}
        \label{fig:markery-stan-polaczenia}
    \end{subfigure}
    \caption{Kodowanie kolorystyczne markerów}
    \label{fig:markery-kolory}
\end{figure}

Węzeł lokalny - czyli ten połączony z telefonem użytkownika bezpośrednio przez Bluetooth -
posiada małą ikonkę telefonu w prawym górnym rogu. Pozostałe węzły mają w tym miejscu kropkę,
której kolor odpowiada jakości połączenia z siecią (zakres SNR\footnote{SNR (ang.\
\textit{Signal-to-Noise Ratio}) - stosunek mocy użytecznego sygnału do mocy szumu,
wyrażony w~dB. Im wyższa wartość, tym lepsze połączenie.} wyrażony w~dB). Strzałka wewnątrz markera wskazuje kierunek
poruszania się węzła (rysunek~\ref{fig:markery-szczegoly}).

\begin{figure}[!h]
    \centering
    \begin{subfigure}{0.48\linewidth}
        \centering
        \includegraphics[width=\linewidth]{markeryJakoscPolaczenia.png}
        \caption{Jakość połączenia (SNR) a kolor kropki}
        \label{fig:markery-jakosc-polaczenia}
    \end{subfigure}
    \hfill
    \begin{subfigure}{0.48\linewidth}
        \centering
        
\includegraphics[width=\linewidth]{markeryKierunek.png}
        \caption{Kierunek poruszania się węzła}
        \label{fig:markery-kierunek}
    \end{subfigure}
    \caption{Wizualizacji markerów}
    \label{fig:markery-szczegoly}
\end{figure}

\paragraph{Geofencing.}

Menu stref dostępne jest pod \textit{floating button} w lewym dolnym rogu ekranu mapy.
Strefy składają się z wielokątów, których wierzchołki tworzone są poprzez długie przyciśnięcie
(\textit{long press}) wybranego punktu na mapie. Po dodaniu trzeciego wierzchołka pojawia się
opcja \textbf{Zamknij}, która tworzy strefę z wprowadzonych punktów
(rysunek~\ref{fig:geofencing-tworzenie}).

\begin{figure}[!h]
    \centering
    \begin{subfigure}{0.45\linewidth}
        \centering
        \includegraphics[width=\linewidth]{strefyGeofencing.jpg}
        \caption{Menu geofencingu - widok pusty}
        \label{fig:geofencing-puste}
    \end{subfigure}
    \hfill
    \begin{subfigure}{0.45\linewidth}
        \centering
        \includegraphics[width=\linewidth]{strefyGeofencingWierzcholki.jpg}
        \caption{Definiowanie wierzchołków strefy}
        \label{fig:geofencing-wierzcholki}
    \end{subfigure}
    \caption{Tworzenie strefy geofencingu}
    \label{fig:geofencing-tworzenie}
\end{figure}

Następnie użytkownik wprowadza nazwę strefy, nadaje jej kolor oraz przypisuje węzły, które
będą wyzwalać zdarzenia wejścia i wyjścia. Dostępne strefy mogą być aktywowane i usuwane
z poziomu listy (rysunek~\ref{fig:geofencing-zarzadzanie}).

\begin{figure}[!h]
    \centering
    \begin{subfigure}{0.45\linewidth}
        \centering
        \includegraphics[width=\linewidth]{nowaStrefa.jpg}
        \caption{Formularz tworzenia nowej strefy}
        \label{fig:nowa-strefa}
    \end{subfigure}
    \hfill
    \begin{subfigure}{0.45\linewidth}
        \centering
        \includegraphics[width=\linewidth]{strefyLista.jpg}
        \caption{Lista zdefiniowanych stref}
        \label{fig:strefy-lista}
    \end{subfigure}
    \caption{Zarządzanie strefami geofencingu}
    \label{fig:geofencing-zarzadzanie}
\end{figure}

\paragraph{Węzły.}

Ekran listy węzłów (\texttt{NodeListScreen}) prezentuje wszystkie aktywne węzły wykryte
w sieci Meshtastic. Dla każdego węzła wyświetlane są jego nazwa, rola, czas ostatniej
aktywności oraz wskaźnik jakości sygnału (SNR) i RSSI\footnote{RSSI (ang.\
\textit{Received Signal Strength Indicator}) - bezwzględna moc odebranego sygnału,
wyrażona w~dBm. W odróżnieniu od SNR nie uwzględnia poziomu szumu, lecz informuje
o sile sygnału w punkcie odbioru.}. Lista może być filtrowana i sortowana,
co ułatwia odnalezienie konkretnego urządzenia w większej sieci. Dotknięcie węzła
otwiera ekran szczegółów (\texttt{NodeDetailScreen}), prezentujący pełne dane
telemetryczne - w tym poziom naładowania baterii, dokładne współrzędne GPS
oraz historię pozycji z bieżącej sesji w uproszczonej formie graficznej (rysunek~\ref{fig:wezly-lista}).

\paragraph{Ustawienia.}

Ekran ustawień (\texttt{SettingsScreen}) umożliwia konfigurację połączenia z aplikacją
Meshtastic. Użytkownik może tu wybrać urządzenie Bluetooth, z którym aplikacja ma się
połączyć, oraz monitorować aktualny stan połączenia. Ekran jest celowo minimalistyczny -
zgodnie z założeniem projektowym, konfiguracja powinna być jednorazowa i nie wymagać
powracania do ustawień podczas normalnej eksploatacji systemu.
